\documentclass[12pt, a4paper]{article}

% --- Pacchetti Scientifici e Formattazione ---
\usepackage{amsmath, amssymb, amsfonts, amsthm}
\usepackage{booktabs, graphicx}
\usepackage[document]{ragged2e}
\usepackage{xcolor}

% --- Codifica e Font ---
\usepackage[utf8]{inputenc}
\usepackage[T1]{fontenc}
\usepackage[italian]{babel}
\usepackage{mathptmx}

% --- Gestione Margini e Layout ---
\usepackage[includeheadfoot]{geometry}
\geometry{
    a4paper,
    top=3cm,
    bottom=3cm,
    left=3.5cm,
    right=3.5cm,
}

% --- Gestione Interlinea e Figure ---
\usepackage{setspace}
\setstretch{1.3}
\usepackage[figuresright]{rotating}
\usepackage[pdftex, hidelinks]{hyperref}

% Document metadata
\title{frontend del configuratore commerciale}
\author{Mattia Cavina}
\date{18 giugno 2026}

\begin{document}

\section{Intrduzione}
Il front end verrà realizzato mediante framework .NET MAUI nativo windows e implementato in 'C\#'.
La scelta di questo framework è stata dettata dalla necessità di avere un'applicazione integrata con i sistemi Windows,
 che possa sfruttare al meglio le funzionalità del sistema operativo e garantire una migliore esperienza utente.
La seconda motivazione è la possibilità di sviluppare un'applicazione cross-platform,
che possa essere eseguita anche su altri sistemi operativi come macOS e Linux o Android,
garantendo una maggiore flessibilità e portabilità del software.

\section{Scelte tecnologiche}
\subsection{Pattern scelto}
Per l'implementazione del frontend mediante .NET MAUI è stato scelto il pattern MVVM (Model-View-ViewModel).
Questo pattern consente di separare la logica di presentazione dalla logica di business,
facilitando la manutenzione e il test del codice.
La View rappresenta l'interfaccia utente, il ViewModel gestisce la logica di presentazione e il Model
rappresenta i dati e le regole di business.

\subsection{Framework classi View}
L'ambiente di sviluppo .NET MAUI mette a disposizione più strade di implementazione delle classi View.
La prima è quella di utilizzare le classi XAML, che consentono di definire l'interfaccia utente in modo dichiarativo tramite 
un linguaggio di markup simile a HTML.
La seconda è quella di utilizzare le classi C\# con framework ".NET Multi-platform App UI (.NET MAUI) Community Toolkit",
che consentono di definire l'interfaccia utente in modo programmatico.
La scelta di utilizzare le classi C\# con il toolkit .NET MAUI Community Toolkit
è stata dettata dalla necessità di avere un maggiore controllo sull'interfaccia utente
e sulla logica di presentazione riducendo la complessità della gestione dell'interfaccia.

\section{Definizione delle interfacce utente}
Nella realizzazione del configuratore la scelta delle interfacce utente è il passo
successivo alla realizzazione dei componenti di getsione dei dati e delle regole di business.
\subsection{Interfacce utente}
Seguendo le inteviste e il precedente sofware sono emerse le seguenti interfacce utente:
\begin{itemize}
      \item Login
      \item Menu principale
      \item Schermata ricerca documenti
      \item Schermata dettaglio documenti
      \item Schermata creazione documenti \& modifica (\ Valutare se separare in tre schermate offerta, ordine, tecnica).
      \item Schermata ricerca clienti
      \item Schermata dettaglio clienti
      \item Schermata creazione clienti \& modifica (Solo campi non importati da ERP o Fusion)
      \item Schermata ricerca elenco prodotti
      \item Schermata dettaglio prodotti \& modifica (Solo campi non importati da ERP o Fusion)
      \item Schermata creazione prodotti commerciali (Solo campi non importati da ERP o Fusion)
      \item Schermata selezione prodotti commerciali da aggiungere al documento
      \item Schermata creazione distinte prodotti commerciali (Sono i MACROGRUPPI di prodotto) \& modifica.
      \item Schermata ricerca utenti
      \item Schermata dettaglio utenti
      \item Schermata creazione utenti \& modifica
      \item Schermata gestione notifiche e feedback
      \item Schermata gestione profilo utente e modifica password
  \end{itemize}
  \subsection{Realizzazione delle prime intefacce}
  Definito l'elenco delle intefacce si è proceduti con il disegnarle almeno nella loro forma base
  madiante Figma.
  Dopo averle disegnate si è fatto un focus group con i commerciali di riferimento
  per raccogliere feedback e suggerimenti.
  \subsection{Prototipo interattivo}
      Dopo aver raccolto i feedback si è proceduto con la realizzazione di un prototipo interattivo
      che permettesse di simulare il flusso di lavoro e le interazioni tra le varie schermate.
      Il prototipo è stato realizzato direttamente in .NET MAUI e permette di navigare tra le varie schermate
      e di testare lo spostamento tra le schermate e la visualizzazione dei dati già presenti nel database.

\end{document}